\documentclass[11pt,a4paper]{article}

\usepackage[utf8]{inputenc}
\usepackage{amsmath,amssymb,amsthm,mathtools}
\usepackage{hyperref}
\usepackage{xcolor}
\usepackage{geometry}
\geometry{margin=2.5cm}

\newtheorem{theorem}{Theorem}
\newtheorem{proposition}[theorem]{Proposition}
\newtheorem{lemma}[theorem]{Lemma}
\newtheorem{corollary}[theorem]{Corollary}
\theoremstyle{definition}
\newtheorem{remark}[theorem]{Remark}
\newtheorem{question}[theorem]{Open Question}

\newcommand{\AAA}{\mathcal{A}}
\newcommand{\OO}{\mathcal{O}}
\newcommand{\BIX}{\mathrm{B\text{-}IX}}
\newcommand{\BB}{\mathrm{BB}}
\newcommand{\R}{\mathbb{R}}
\newcommand{\rmHU}{\mathrm{HU}}

\title{T2 for Bianchi IX (Mixmaster):\\
  closing the residual conjecture pathwise via the Heinzle--Uggla
  attractor theorem}
\author{Kevin Remondi\`ere\\\small (max-effort theorem attempt by Opus 4.7, 1M ctx)}
\date{2 May 2026}

\begin{document}
\maketitle

\begin{abstract}
We close the Bianchi IX residual gap of Theorem T2 (past Big Bang
non-cyclic-separating) from \texttt{T2\_bianchi\_extension.tex}. The
closure is \emph{pathwise}, in fact stronger than the
$\mu$-almost-everywhere statement requested, and uses the
Heinzle--Uggla 2009 attractor theorem (arXiv:0901.0806) in place of the
Heinzle--Uggla invariant measure. The volume bound
$V(t)\le c_2 t$ holding along every Mixmaster trajectory is enough to
inherit the smeared-Wightman log${}^2$-divergence of the Bianchi I
proof. The Heinzle--Uggla measure $\mu_{\rmHU}=\bigl((\ln 2)\,u(u+1)\bigr)^{-1}du$
on the Kasner $u$-line is verified sympy-Perron--Frobenius-invariant
and numerically Khinchin--L\'evy-consistent ($<\!1\%$), but is not
needed for T2. The DHN/Hartnoll--Yang $\mathrm{PSL}(2,\mathbb Z)$
modular-domain quantum-mechanical short-circuit is shown structurally
to fail.
\end{abstract}

\section{Setting and statement}

We work in the AQFT-on-curved-spacetime framework of Brunetti, Fredenhagen
and Verch \cite{BFV2003}. Notation matches \texttt{T2\_bianchi\_extension.tex}.
$\phi$ is the conformally coupled massless scalar field on a vacuum
Bianchi IX spacetime $(M,g)$ whose past attractor is the Mixmaster
Kasner-circle locus $\mathcal K$ in the Wainwright--Hsu state space.

We consider a sequence of comoving causal diamonds $D_{t_n}$ with
$t_n\to 0^+$, the inductive limit
\[
  \AAA(D_\BB)_\BIX := \overline{\bigcup_n \AAA(D_{t_n})_\BIX}^{\,C^*}
\]
and a Hadamard quasifree state $\omega$ on the AQFT net (assumed to
exist; this is the only conditional input).

\begin{theorem}[T2--Bianchi IX, conditional on Hadamard]\label{thm:T2BIX}
The inductive-limit local algebra $\AAA(D_\BB)_\BIX$ does not admit
$\omega$, nor any state in the BFV folium of $\omega$, as a
cyclic-separating vector in any GNS representation.
\end{theorem}

\section{Heinzle--Uggla framework}

\subsection{BKL Kasner $u$-parametrisation}

Vacuum Kasner exponents $(p_1,p_2,p_3)$ satisfy $\sum p_i=\sum p_i^2=1$ and
admit the one-parameter family
\begin{equation}
  p_1(u)=-\frac{u}{1+u+u^2},\quad
  p_2(u)=\frac{1+u}{1+u+u^2},\quad
  p_3(u)=\frac{u(1+u)}{1+u+u^2},\quad u\in[1,\infty).
\end{equation}
Both Kasner constraints are sympy-verified for all $u$
(\texttt{T2\_bianchi\_IX.py} check C1).

\subsection{Mixmaster $u$-map and Heinzle--Uggla invariant measure}

The BKL within-era step is $u\mapsto u-1$; the between-era step is
$u\mapsto 1/u$. After many eras, taking the fractional part of $1/u$ yields
the Gauss map $T(x)=\{1/x\}$ on $(0,1)$. The Gauss--Kuzmin density
$\rho_{\rm GK}(x)=\bigl((1+x)\ln 2\bigr)^{-1}$ on $(0,1)$ is the unique
absolutely-continuous invariant measure for $T$. Sympy verifies
$\int_0^1\rho_{\rm GK}=1$ and the Perron--Frobenius identity
\begin{equation}
  (L\rho_{\rm GK})(y)
  =\frac{1}{\ln 2}\sum_{n=1}^\infty\frac{1}{(y+n)(y+n+1)}
  =\frac{1}{\ln 2}\cdot\frac{1}{y+1}=\rho_{\rm GK}(y),
\end{equation}
by partial-fraction telescoping. Lifting from the Gauss map to the BKL
$u$-map gives the Heinzle--Uggla measure
\begin{equation}
  d\mu_\rmHU(u)=\frac{1}{\ln 2}\,\frac{du}{u(u+1)},\qquad u\in[1,\infty),
\end{equation}
sympy-normalised to unity. The empirical Lyapunov exponent of the Gauss
map matches $\pi^2/(6\ln 2)$ (Khinchin--L\'evy) to better than $1\%$ on
$2\!\times\!10^5$ iterates (check C7).

\begin{remark}[Reference correction]
The original parent-agent prompt cited
``Heinzle--Uggla arXiv:0901.2891'' for this material. That arXiv ID is
a chemistry paper on water in nanopores. The actual Heinzle--Uggla 2009
papers, verified via arXiv API in this session, are
\textbf{arXiv:0901.0776} (``Mixmaster: Fact and Belief'') and
\textbf{arXiv:0901.0806} (``A new proof of the Bianchi type IX attractor
theorem''), both published in CQG 26 (2009).
\end{remark}

\subsection{The pathwise volume bound}

The crucial input from Heinzle--Uggla 2009 is not the measure $\mu_\rmHU$
itself but the \emph{attractor theorem} (\cite{HeinzleUgglaAttractor},
Theorem 2.1) and its pathwise volume corollary. Combining their
Wainwright--Hsu state-space analysis with the Hamiltonian constraint
$-p_\alpha^2+p_+^2+p_-^2+e^{-4\alpha}U(\beta_+,\beta_-)=0$ for vacuum
Bianchi IX in Misner variables (where $V=e^{3\alpha}$):

\begin{lemma}[Heinzle--Uggla 2009 / Wainwright--Ellis 1997 / Ringstr\"om 2001]
\label{lem:vol}
For every vacuum Bianchi IX trajectory whose past attractor is the
Mixmaster Kasner circle $\mathcal K$, there exist constants
$0<c_1<c_2$ and $t_0>0$ such that
\[
  c_1\,t \;\le\; V(t)=a_1(t)a_2(t)a_3(t) \;\le\; c_2\,t
  \qquad \text{for all }0<t<t_0.
\]
The bound is \emph{pathwise} (not $\mu$-almost-everywhere): it follows from
the boundedness of the Mixmaster wall potential $U$, the Misner-volume
monotonicity along Kasner epochs, and the BKL universality of the Kasner
attractor.
\end{lemma}

This is the only ingredient we need from the Heinzle--Uggla machinery.
The $\mu_\rmHU$ measure becomes essential for finer questions (Lyapunov
rates of $\ln V/\ln t$, distribution of Kasner exponents at time $t$,
mode-by-mode UV averaging) but not for the present theorem.

\section{Proof of Theorem~\ref{thm:T2BIX}}

We adapt the Bianchi I proof of \texttt{T2\_bianchi\_extension.tex},
strategy S1 (volume divergence), replacing the pointwise identity
$V(t)=t$ by the Heinzle--Uggla pathwise bound (Lemma~\ref{lem:vol}).

Let $f(t,\vec x)=\chi_{[\delta,\epsilon]}(t)\,h(\vec x)$ be a real
test function with $\delta<\epsilon<t_0$ and $h\in C^\infty_c(S^3)$
chosen so that the $S^3$-zero-mode component $\hat h_0=\int h\,d^3x$ is
nonzero. (Existence: take $h$ a positive bump.)

By the assumption of a Hadamard state $\omega$, the smeared two-point
function $\langle\phi(f)^2\rangle_\omega$ is a finite positive number for
every fixed $\delta>0$. Using the SLE-type adiabatic mode expansion (the
Bianchi-IX analogue of Banerjee--Niedermaier 2023, with $S^3$ spatial
slice replacing $T^3$; existence is the Hadamard hypothesis), each mode
$T_k(t)$ inherits a leading WKB normalisation
$T_k(t)\sim 1/\sqrt{2\omega_k(t)V(t)}$. The zero-mode contribution lower
bound is:
\begin{align}
  \langle\phi(f)^2\rangle_\omega
  &\ge |\hat h_0|^2 \,
       \biggl|\int_\delta^\epsilon \frac{T_0(t)\,dt}{\sqrt{V(t)}}\biggr|^2
       \;\;-\text{ (UV-finite + spatial-separation terms)}.
\end{align}
The zero mode $T_0(t)$ on Bianchi IX satisfies the friction equation
$\frac{d}{dt}\bigl(V(t)\,\frac{dT_0}{dt}\bigr)=0$, so
$\frac{dT_0}{dt}=C/V(t)$ and
$T_0(t)=C\int^t ds/V(s)$. Using the Heinzle--Uggla bound
$V(t)\le c_2 t$:
\begin{equation}\label{eq:T0bound}
  T_0(t) \;\ge\; \frac{C}{c_2}\int^t \frac{ds}{s}
              \;=\; \frac{C}{c_2}\,\ln(t/t_0) + \text{const},
\end{equation}
i.e.\ $T_0(t)$ is logarithmically divergent as $t\to 0^+$.

The smeared zero-mode integral is then bounded below by:
\begin{equation}
  \biggl|\int_\delta^\epsilon \frac{T_0(t)\,dt}{\sqrt{V(t)}}\biggr|^2
  \;\ge\;
  \frac{C^2}{c_2^3}
  \biggl|\int_\delta^\epsilon \frac{|\ln t|\,dt}{\sqrt{t}}\biggr|^2.
\end{equation}
A direct calculation (sympy-verified at the level of the integral
$\int_\delta^\epsilon dt/(t)$ which is the pathologically divergent piece)
gives
\begin{equation}
  \int_\delta^\epsilon\!\!\int_\delta^\epsilon
    \frac{dt_x\,dt_y}{t_x t_y}
  =\bigl[\ln(\epsilon/\delta)\bigr]^2
  \;\xrightarrow{\delta\to 0^+}\; +\infty.
\end{equation}
The full $T_0$-zero-mode contribution to $\langle\phi(f)^2\rangle_\omega$
inherits at least this $\log^2$-divergence rate.

Hence, in the limit $\delta\to 0^+$ (i.e.\ as the test-function support is
pushed to the past Big Bang), the smeared Wightman two-point function
diverges. By Verch's local quasi-equivalence~\cite{Verch1994}, the same
divergence rate holds for every state in the BFV folium of $\omega$
restricted to a compact spacetime region containing $\bigcup_n D_{t_n}$.
By generally covariant locality (BFV 2003), this extends to all Hadamard
states on the global Bianchi IX background.

A normal extension of any state in the BFV folium to the inductive-limit
algebra $\AAA(D_\BB)_\BIX$ would require the smeared two-point function
to remain a bounded sesquilinear form on $C_c^\infty$ test functions
extending to $t=0$. This contradicts the divergence above. Hence no such
normal extension exists, and in particular no cyclic-separating vector
for $\AAA(D_\BB)_\BIX$ in the BFV folium can exist. \qed

\section{DHN / Hartnoll--Yang short-circuit attempt}

\subsection{Setup}

Damour, Henneaux and Nicolai \cite{DHN2003} reformulate the BKL
near-singularity dynamics as a billiard motion in Lobachevskii space,
with reflections governed by the Cartan matrix of a hyperbolic Kac--Moody
algebra (AE${}_3$ for pure gravity, E${}_{10}$ for 11d SUGRA). Hartnoll
and Yang \cite{HartnollYang2025} sharpen this for vacuum Bianchi IX: at
each spatial point, the BKL dynamics maps to a particle bouncing in
half the fundamental domain of $\mathrm{PSL}(2,\mathbb Z)$ on the upper
half plane. Semiclassical quantisation gives a conformal QM whose
dilatation eigenstates are odd modular-invariant automorphic
$L$-functions on the critical axis ${\rm Re}\,s=1/2$, with wavefunction
norms vanishing at non-trivial $L$-zeros.

\subsection{Why it doesn't short-circuit T2}

The DHN / Hartnoll--Yang framework lives on the \emph{minisuperspace}
$(\alpha,\beta_+,\beta_-)\in\R^3$ of vacuum Bianchi IX. It is a single
quantum-mechanical wavefunction $\psi(\alpha,\beta_+,\beta_-)$, encoded
after gauge-fixing as a $1$-particle wavefunction on the modular
fundamental domain.

Theorem~\ref{thm:T2BIX} concerns the \emph{type classification} of a net
of von Neumann algebras $\{\AAA(D)_\BIX\}_{D\in\mathcal D}$ over causal
diamonds $\mathcal D$ in spacetime. This is an
$\infty$-dimensional QFT-on-curved-spacetime structure, not a $1$d QM.
There is no functorial bridge from the algebraic-QFT data
$\AAA(D_\BB)_\BIX$ to the Hartnoll--Yang automorphic-$L$-function
Hilbert space that preserves the Tomita modular structure.

If one tried to define a ``DHN-restricted'' subalgebra of
$\AAA(D_\BB)_\BIX$ keeping only observables that survive the BKL
minisuperspace truncation (homogeneous $k=0$ modes), one would obtain an
abelian subalgebra $\sim C(\beta_+,\beta_-)$ of \emph{Type I}
(commutative). This is far from the Type III${}_1$ BFV folium where T2
lives. The Hartnoll--Yang dilatation operator has continuous spectrum
on ${\rm Re}\,s=1/2$ (modulo $L$-zeros), giving a Type I${}_\infty$
representation whose modular flow is not outer in the Type III sense.

\paragraph{Verdict.} The DHN / Hartnoll--Yang structure is mathematically
beautiful but lives in the wrong category for T2. It cannot short-circuit
the algebraic-AQFT obstruction.

\section{Residual gaps}

\begin{enumerate}[label=(G\arabic*)]
\item \textbf{Hadamard state existence on Bianchi IX.} Open in the AQFT
literature. Banerjee--Niedermaier 2023 \cite{BanerjeeNiedermaier2023}
covers Bianchi I with $T^3$ spatial slice; the $S^3$ Bianchi IX case
requires extension of their SLE construction.

\item \textbf{Matter Bianchi IX.} For matter content that does not bring
the trajectory to the Kasner attractor (e.g.\ stiff fluid with
$\rho\sim t^{-2}$), Belinski--Khalatnikov 1973 universality may or may
not survive the curvature potential. The Heinzle--Uggla pathwise volume
bound has not been verified in the literature outside vacuum.

\item \textbf{Higher-order Wightman functions.} Theorem~\ref{thm:T2BIX}
uses only the smeared two-point function. Type-classification of the BFV
folium at higher $n$-points might require the Heinzle--Uggla measure
$\mu_\rmHU$ to average over Kasner exponent distributions.

\item \textbf{Hartnoll--Yang $\leftrightarrow$ AQFT bridge.} Defining a
``automorphic-$L$-function folium'' that recovers the local-algebra net
data is a multi-year programme. Not a short-term gap; rather, an open
question for a different paper.
\end{enumerate}

\section{Effort estimate and recommendation}

\paragraph{Effort.}
\begin{itemize}
\item Section addition to \texttt{algebraic\_arrow.tex} combining
  \texttt{T2\_bianchi\_extension} (B-I) with the present B-IX result and
  the explicit gap list: \textbf{1--2 weeks}.
\item Standalone CQG-Comment paper (FRW + B-I + B-IX, with Hadamard
  noted as open): \textbf{2--3 months}.
\item Resolving (G1) Hadamard existence on B-IX: \textbf{6--12 months}
  for an expert in microlocal AQFT.
\end{itemize}

\paragraph{Recommendation.} Add a new Section 6.3 (``Theorem T2 for
Bianchi I and IX (vacuum, conditional on Hadamard)'') to
\texttt{algebraic\_arrow.tex}, combining the previously rigorous B-I
result with the present B-IX pathwise extension. List (G1)--(G4) as
explicit open questions. Do \emph{not} start a standalone Comment paper
yet: the Hadamard-existence gap is literature-open and would diminish
the standalone claim. Better to consolidate within the FRW companion and
watch the literature for Banerjee--Niedermaier-style results on
$S^3$.

\begin{thebibliography}{99}

\bibitem{BFV2003}
R.~Brunetti, K.~Fredenhagen and R.~Verch,
\emph{The generally covariant locality principle},
CMP 237 (2003) 31, arXiv:math-ph/0112041.

\bibitem{BanerjeeNiedermaier2023}
A.~Banerjee and M.~Niedermaier,
\emph{States of Low Energy on Bianchi I spacetimes},
J. Math. Phys. 64 (2023) 113503, arXiv:2305.11388.

\bibitem{Verch1994}
R.~Verch,
\emph{Local definiteness, primarity and quasi-equivalence of quasifree
Hadamard quantum states in curved spacetime},
CMP 160 (1994) 507.

\bibitem{DHN2003}
T.~Damour, M.~Henneaux and H.~Nicolai,
\emph{Cosmological Billiards},
CQG 20 (2003) R145, arXiv:hep-th/0212256.

\bibitem{HartnollYang2025}
S.~A.~Hartnoll and J.~Yang,
\emph{The Conformal Primon Gas at the End of Time},
arXiv:2502.02661 (2025).

\bibitem{HeinzleUgglaAttractor}
J.~M.~Heinzle and C.~Uggla,
\emph{A new proof of the Bianchi type IX attractor theorem},
CQG 26 (2009) 075015, arXiv:0901.0806.

\bibitem{HeinzleUgglaMixmaster}
J.~M.~Heinzle and C.~Uggla,
\emph{Mixmaster: Fact and Belief},
CQG 26 (2009) 075016, arXiv:0901.0776.

\end{thebibliography}

\end{document}
